\pdfoutput=1
\documentclass[11pt]{article}
\usepackage[T1]{fontenc}
\usepackage{lmodern,microtype}
\usepackage[hmargin=1in,vmargin=0.9in]{geometry}
\usepackage{amsmath,amssymb,amsthm,booktabs}
\usepackage[hidelinks]{hyperref}
\hypersetup{pdftitle={γ(Q₄₇) = 25 and bounds for queen domination on odd boards},pdfauthor={Jason Davies}}
\newtheorem{theorem}{Theorem}[section]
\newtheorem{corollary}[theorem]{Corollary}
\newtheorem{lemma}[theorem]{Lemma}
\newtheorem{proposition}[theorem]{Proposition}
\newcommand{\ind}[1]{\mathbf{1}_{#1}}
\newcommand{\file}[1]{\nolinkurl{#1}}
\title{$\gamma(Q_{47})=25$ and bounds for queen domination on odd boards}
\author{Jason Davies}
\date{27 November 2026}
\begin{document}
\maketitle

\begin{abstract}
The domination and independent domination numbers of the $47\times47$
queen's graph are both 25. The independent upper bound is due to
\"Osterg{\aa}rd and Weakley. We exhibit a placement and exclude every placement of 24 queens by a
computer-assisted proof. A geometric bound eliminates 271,298 of 369,412
cases on canonical frame rectangles. The remaining 98,114 cases have
91,017 representatives under board symmetry. We close them by SAT
refutations and arithmetic exclusions. In the final 15 branches the occupied rows and
columns are forced. The targets on empty lines constrain the diagonal
values, and their first and second moments rule out every remaining
combination. We also prove that on a board of order $4m-1$, a nonorthodox
dominating set of $2m$ queens has at least $m-1$ queens of each colour.
\end{abstract}

\section{The domination number at order 47}\label{sec:bounds}
The vertices of $Q_n$ are the squares of the $n\times n$ board. Queens
dominate the board if every square contains a queen or shares a row,
column or diagonal with one. Write $\gamma(Q_n)$ for the minimum number
of queens with this property, and $i(Q_n)$ for the minimum size of a
dominating set whose queens are pairwise nonattacking.
\begin{theorem}\label{thm:q47}
\(\gamma(Q_{47})=i(Q_{47})=25\).
\end{theorem}
The lower bound for odd boards gives
$\gamma(Q_{47})\ge24$ \cite{FW}. For the other direction, number the rows
and columns from 0 to 46, writing each square as (column, row). Place a queen in column $2i$ and row $r_i$ for
$0\le i<24$, where the successive rows are
\[
\begin{array}{rrrrrrrrrrrr}
15&9&39&33&11&41&17&9&25&37&43&23\\
7&1&21&27&3&45&35&29&19&5&31&13 .
\end{array}
\]
Testing the four lines through each square leaves just $(7,46)$ uncovered.
Adding a queen there proves $\gamma(Q_{47})\le25$. For independent domination,
\"Osterg{\aa}rd and Weakley \cite[pp.~10--11]{OW} exhibit an independent
dominating 25-set. The verifier adds 23 to both of their centered coordinates
and checks that the set occupies 25 distinct lines in each family and
dominates all 2,209 squares. Thus $25=\gamma(Q_{47})\le i(Q_{47})\le25$
once the lower bound is established.

The frame method of \cite{Davies}, following the box decomposition of
Rostami~\cite{HC}, resolved the odd board $Q_{35}$. Its even-board applications
have since been extended to $Q_{48}$ by \DJ uki\'c \cite{Djukic} and to the
even boards from 50 through 64 by Chan \cite{Chan}. The latter work counts
perimeter incidences separately by square colour. We adapt that idea in
Section~\ref{sec:colour}. At order 47 the geometry allows more unused
incidences than in the even-board problem with $n/2$ queens. Keeping track
of where they can occur is useful both before and during SAT solving.

\section{An odd-board reduction before SAT}
Let $n=2k-1\ge5$ and suppose $k$ queens dominate $Q_n$. The leftmost and
rightmost empty columns $a,b$ and the lowest and highest empty rows $c,d$
bound a rectangle called the \emph{frame}. Its spans are $w=b-a$ and
$h=d-c$. For each of the four families of lines, subtract the number of
occupied lines from $k$ to obtain deficits $\delta_R,\delta_C,\delta_T,
\delta_S$ (row, column, difference and sum). An empty row and an empty
column meet in a square requiring a diagonal attack. Call it a
\emph{target}. There are $k-1+\delta_R$ empty rows and $k-1+\delta_C$
empty columns.

\begin{proposition}\label{prop:margin}
Let $M=a+(n-1-b)+c+(n-1-d)$ and let $L$ count queens $(x,y)$
with $x\notin[a,b]$ and $y\notin[c,d]$. Then
\[
 M-k\le L\le4-\delta_R-\delta_C.
\]
When $w\ge h$ we also have $L\le2-2\delta_R+2\ind{w=h}$.
When $h\ge w$ the corresponding bound uses $\delta_C$.
\end{proposition}
\begin{proof}
All $M$ exterior rows and columns are occupied. A queen contributes at
most one incidence with these lines except for the $L$ queens counted
twice. Each of those $L$ queens has a diagonal missing the frame. The
$2k$ diagonal occurrences therefore cover at most $2(2k-L)$ perimeter
targets. Counting targets on the four sides and subtracting the four
corners gives $2(k-1+\delta_R)+2(k-1+\delta_C)-4$. Comparing these counts
gives the first two inequalities. When $w\ge h$, an occupied diagonal
can meet the vertical sides in at most one point, apart from the two
corner-to-corner diagonals when $w=h$. The two sides have
$2(k-1+\delta_R)$ targets. They receive at most $2k-L+2\ind{w=h}$
distinct hits. Interchanging the axes finishes the proof.
\end{proof}
This refines the exterior-queen count in the rectangular board argument
of Boz\'oki, G\'al, Marosi and Weakley \cite{BGMW}.

\begin{lemma}\label{lem:grid}
Let $A,B$ be arbitrary nonempty finite sets of coordinates. Suppose the diagonals of $s$ queens cover
$A\times B$. Then $|A|+|B|\le2s+2$. At equality both sizes are even.
\end{lemma}
\begin{proof}
If one size is one, at most $2s$ points on that line can be hit. Otherwise
the border of the grid has $2|A|+2|B|-4$ points, and each of the $2s$
diagonals hits at most two. At equality $|A|+|B|=2s+2$ is even. It remains to rule out two odd
sizes. We use the coefficient argument of Alon's Combinatorial
Nullstellensatz \cite[Theorem~1.2]{Alon}. Consider
\[
 P(x,y)=\prod_{(u,v)}\big((x-u)^2-(y-v)^2\big),
\]
with a factor for each queen. It vanishes on the grid. Its remainder
modulo $\prod_{a\in A}(x-a)$ and $\prod_{b\in B}(y-b)$ is identically
zero by univariate interpolation twice. At equality the monomial
$x^{|A|-1}y^{|B|-1}$ has degree $2s$ and is already reduced. Its
coefficient survives reduction, because reducing any other monomial
strictly lowers its degree. Both exponents are even, so its coefficient
in $(x^2-y^2)^s$ is a nonzero signed binomial coefficient.
\end{proof}

Weakley \cite{Weakley} proved $\gamma(Q_{4k+1})\ge2k+1$.
Burger and Mynhardt \cite{BM} studied dominating sets on $Q_{4k+3}$,
the congruence class considered here. We distinguish placements by their
occupied coordinate parities and count queens of each square colour.

Call a cover \emph{orthodox} if it occupies all rows of some parity and all
columns of some parity, following Neuhaus \cite{Neuhaus}. Colour squares by
$x+y\bmod2$. Orthodoxy concerns occupied lines, while square colour
concerns the queens themselves.

For even $n\ge14$, Neuhaus \cite[Lemma~4.4]{Neuhaus} proved that in a
nonorthodox dominating set of $n/2$ queens, each colour class has size
between $n/4-1/2$ and $n/4+1/2$. His proof counts diagonal coverage on
products of empty lines of specified parity. The following odd-board
analogue uses the obstruction to odd--odd equality in Lemma~\ref{lem:grid}.

\begin{theorem}\label{thm:split}
On $Q_{4m-1}$, $m\ge2$, the smaller colour class in a dominating set of
$2m$ queens has size $0$ or $1$ in the orthodox case, and $m-1$ or $m$
in the nonorthodox case. A nonsquare frame with odd longer span forces
exactly $m$ queens of each colour.
\end{theorem}
\begin{proof}
Put $k=2m$. In an orthodox cover at most one queen can miss the compulsory
parity on each axis, since the two parity classes have sizes $k$ and
$k-1$. Two queens of the opposite colour would have to be distinct
exceptions, one on each axis. There would then be $k-1$ empty indices of
the opposite parity on each axis. Their product has the compulsory square
colour and must be covered by the diagonals of the other $k-2$ queens.
This attains equality in Lemma~\ref{lem:grid} with two odd sizes.

For a nonorthodox cover, suppose one axis has empty lines of only one
parity. The other has empty lines of both parities. Two opposite-parity
empty lines on the latter axis each contain at least $k-1$ targets, of
one colour on the first and the other colour on the second. Each queen
hits at most two targets on either line, forcing at least $m$ of each
colour. Otherwise all four empty parity classes are nonempty. For a
given square colour the targets consist of two products of these classes.
Adding Lemma~\ref{lem:grid} over the two products counts each empty
index once and yields $4m-2+\delta_R+\delta_C\le4s+4$, where $s$ is the
number of queens of this colour. Thus $s\ge m-1$.

Finally, orient the longer span horizontally. If it is odd, every empty
row meets the two extreme columns in one target of each colour. A diagonal
can reach at most one of the two vertical sides. Each colour therefore
needs at least $\lceil(k-1+\delta_R)/2\rceil\ge m$ queens.
\end{proof}

At $n=47$ the two possible nonorthodox splits are $11+13$ and $12+12$.
Neuhaus reports an integer-programming exclusion of orthodox 24-queen
covers \cite[\S3.6, p.~59]{Neuhaus}. Thus a nonorthodox solution would
require moving at least eleven of the monochromatic queens displayed in
Section~\ref{sec:bounds}. Our exhaustive search includes the orthodox
cases and verifies their exclusion independently of that result.

\section{Square frames and colour-specific perimeter counts}\label{sec:colour}
For a square frame translate its lower left corner to the origin and write
$w$ for its span. Pair each interior offset $i$ with $w-i$, counting the
middle offset only once. For each resulting set $J$, let $r_J,c_J$ count
its empty rows and empty columns. Put $E=\sum_J|r_J-c_J|$.

\begin{proposition}\label{prop:profile}
We have $E\le4-\delta_R-\delta_C-\delta_T-\delta_S$. When $w$ is odd,
the number $z$ of queens of either colour satisfies
$2z\ge k-2+(\delta_R+\delta_C+E)/2$.
\end{proposition}
\begin{proof}
For $1\le i<w$, group together the points
$(0,i),(i,0),(w,w-i),(w-i,w)$. A diagonal through any such point can
meet the perimeter only within its group, at most once on a vertical
side and once on a horizontal side. If a group has $v_i,h_i$ vertical
and horizontal targets it needs at least $\max(v_i,h_i)$ distinct
diagonals. The corners need at least two more. Now
\[
 \sum v_i=2(k-3+\delta_R),\qquad
 \sum h_i=2(k-3+\delta_C),\qquad \sum |v_i-h_i|=2E .
\]
For the last identity use the two groups indexed by each reflected pair.
At a middle offset each empty line contributes two targets. Summing
$\max(v_i,h_i)=(v_i+h_i+|v_i-h_i|)/2$ and comparing with the
$2k-\delta_T-\delta_S$ occupied diagonals gives the first bound. When
$w$ is odd the two groups of a reflected pair have opposite colours
and equal target counts. Each colour also needs one diagonal for its
corners. The same count for a single colour gives the second bound.
\end{proof}

We next specialise to a square frame with zero deficits. Each queen then
has a different line in each of the four families. Fix one colour. Let
$z$ be its queen count, $v,h$ its numbers of vertical- and horizontal-side
targets, $c_0$ its number of corners, and $b_0$ the number of occupied
corner-to-corner diagonals of this colour. Define waste as the capacity
for diagonal hits less the number of distinct targets covered. On the full
perimeter each occupied diagonal has capacity two. On opposite sides it
has capacity one, increased to two for a corner-to-corner diagonal.
Missing a side, hitting a nontarget or hitting a target twice can all
use up capacity. In a cover the three wastes are exactly
\[
 W_F=4z-v-h+c_0,\qquad W_V=2z+b_0-v,\qquad W_H=2z+b_0-h .
\]
Indeed the $z$ queens supply $2z$ distinct diagonals. All $v+h-c_0$
perimeter targets must be hit, with each corner counted once. The
side identities follow by counting targets on the two sides in question.

For odd $w$, $v=h=k-1$ and $c_0=2$ for each colour. Thus at order 47 a
colour represented by eleven queens has $W_F=0$. Every one of its 22
diagonals must reach two distinct targets and all 44 targets must be
covered once. For even $w$ all corners have the same colour. For the
other colour $v,h$ are even and $b_0=c_0=0$, so $W_V,W_H$ are even and
$W_F=W_V+W_H$. The total full-perimeter waste across colours is eight.
The total on either pair of sides is two plus the number of occupied
corner-to-corner diagonals. The encoder enumerates the nonnegative integer waste counts satisfying
these identities and bounds and enforces the counts for the selected
budget of each colour.

\section{A shared certificate for an orthodox profile}\label{sec:orthodox}
Consider placements with all 24 even rows and all 23 odd columns occupied
and no queen in any other column. Each even row carries one queen. One
odd column $x_*$ carries two, and every other odd column carries one.
An entire family of frame branches forcing this pattern can use the
same exclusion, also after transposing the board.

For each queen put $T=y-x$ and $S=x+y-46$. The targets form the product
of the even columns with the odd rows. A target $(x,y)$ requires its
$T$ or $S$ to occur among the queen indices. For each target difference
$t$, list the sum indices $N(t)$ compatible with it. Solving $y-x=t$, $x+y-46=s$ for a target gives
$s\equiv t-46\pmod4$, $|s-t|\le46$ and $|s+t|\le44$.
For $t>0$ the interval for $s$ is centered at $-1$, and for $t<0$
at $1$, in both cases with radius $45-|t|$. Within a fixed residue
class of $t$ the admissible $s$ have the same residue. Distinct values
of $|t|$ differ by at least two, which is at least the distance between
the centers. Thus the neighbourhoods are nested in order of increasing
$|t|$.
Before the first absent difference in this order all differences are
compulsory. At that first absence all sums in $N(t)$ are compulsory. An
additional choice represents the case where every difference occurs.
Choose a prefix independently for the two residue classes. Allowing the
first absent difference back into a multiset completion only enlarges the
set of candidates.

Complete each compulsory list to a multiset of 24 indices, allowing repeats.
With $u=23$ and $A=\sum_{i=0}^{46}(i-u)^2$, necessary conditions are
\[
 \sum T=u-x_*,\qquad \sum S=-\sum T,\qquad
 \sum T^2+\sum S^2=2\big(A+(x_*-u)^2\big).
\]
The linear sums follow by adding the known row and column coordinates.
The quadratic identity follows from $T^2+S^2=2[(x-u)^2+(y-u)^2]$.
Moreover $T\equiv S\bmod4$ for a queen in an odd column, so the
multiplicities in residue class one agree. For each prefix pair the
generator enumerates the two completions separately over the odd indices
from $-45$ through $45$, and compares their
linear sum, square sum and residue count. Every placement would yield
one matching pair of keys.

After the list-length test there are 70 prefix choices. The generated
formula has 10,920 variables and 10,921 clauses. At least one prefix must
be selected, and a selected prefix must have a compatible key. The
enumeration rules out every such key. A 563-byte DRAT trace eliminates the
70 selectors and derives the empty clause. Its small size reflects the
work done by the multiset enumerator. Checking the trace establishes
unsatisfiability of this formula. Applying it to a branch also requires
the frame and its fixed line occupancies to imply the stated profile.
The trace checks the generated clauses. The multiset enumeration supplies
their mathematical justification. A second checker enumerates unions
of the difference-neighbourhoods of omitted sums directly on the grid.
Each such union is compulsory. Every sum with a neighbour outside the
union must be present. This produces a subcover of every possible
placement. Completing these subcovers to multisets gives 70 support
choices and no matching moment keys, independently of the prefix order,
CNF generator and DRAT checker.

\begin{proposition}\label{prop:axis}
For a branch, enumerate the occupied row and column counts of each parity
consistent with the layer bits and the prescribed line totals. For each
choice, distribute 24 queens over the four row--column parity classes.
Each occupied line requires at least one queen. A parity class with no
occupied lines carries no queens. If there are $a>0$ empty rows of parity $p$ and $b>0$
empty columns of parity $q$, require at least $\lceil\max(a,b)/2\rceil$
queens of colour $p+q\bmod 2$. If no distribution survives, the branch
is impossible. If all survivors have a profile excluded for the prescribed diagonal
deficits, the corresponding profile proof closes the branch.
\end{proposition}
\begin{proof}
Every placement supplies one of the enumerated distributions. On any fixed
empty line, a queen can hit at most two targets diagonally, and can hit
only targets of its colour. Apply this observation to a row and to a
column of the indicated grid. This proves the required lower bound and
both conclusions.
\end{proof}

For example, suppose $\delta_R=0$ and $\delta_C=1$. If the prescribed
occupancies already force $r$ distinct even rows and $s$ distinct odd
columns, at most $U=(24-r)+(24-s)$ queens can have even square colour:
each such queen misses at least one of these two coordinate parities.
An empty even row would have at least $1+s$ targets in even columns,
since at most $23-s$ even columns are occupied. Hence $1+s>2U$ forces
all even rows to be occupied. Similarly, an empty odd column would meet
at least $r-1$ empty odd rows, so $r-1>2U$ forces all odd columns to
be occupied. Either deduction tightens the count for the other. Once
all odd columns are occupied, the even-colour queens must lie in odd
rows and there are at most $24-r$ of them. These tests certify the
profile, for example, in the case with margins $(0,0,1,1)$, deficits
$(0,1,0,0)$ and the 76-bit cube $(11110000)^9\,1111$ in the layer
ordering below.

\section{Diagonal moments on the final 15 branches}\label{sec:terminal}
In 15 branches from 13 frames, the fixed layer bits and line counts
determine both sets of occupied coordinates. Six branches occupy exactly
the even coordinates on both axes, three occupy exactly the odd coordinates
on both axes, and six occupy all odd coordinates and one even coordinate
on each axis.

For the first two families write each queen as $(p+2i,p+2j)$ with
$p=0$ or $1$. Set $L=23-p$, $D=i-j$ and $H=i+j-L$.
When $p=0$, each axis uses the ranks $0,\ldots,L$ once, so
$\sum D=\sum H=0$. When $p=1$, one column rank $I$ and one row rank
$J$ occur twice. Then $\sum D=I-J$ and $\sum H=I+J-L$.
In both families,
\[
 \sum D^2+\sum H^2-(\sum D)^2-(\sum H)^2
 =\sum_{i=0}^L(2i-L)^2 .
\]
Indeed, $D^2+H^2=((2i-L)^2+(2j-L)^2)/2$. For the odd axes the two
repeated ranks add $((2I-L)^2+(2J-L)^2)/2$, exactly the sum of the
squared first moments.

In the mixed family every occupied line carries one queen. If distinct
queens occupied the two even lines, both queens would have odd square
colour. The $23\times23$ product of empty even lines would have to be
covered by the diagonals of the other 22 queens. This contradicts the
equality case of Lemma~\ref{lem:grid}. The queen on the even column
therefore lies on the even row. Removing it leaves a permutation on the
23 odd coordinates. Its diagonal ranks satisfy $\sum D=\sum H=0$ and
$\sum(D^2+H^2)=\sum_{i=0}^{22}(2i-22)^2=4048$.

Write the rank pair of a target $(x,y)$ as
$(d,h)=((x-y)/2,(x+y-46)/2)$. A target on two empty lines
requires its difference or its sum in the corresponding diagonal multiset.

\begin{lemma}\label{lem:pure-enum}
For either same-parity family every cover occurs among the following
pairs of multiset completions. In each parity class of target differences
choose the least magnitude at which some signed difference is absent,
or the next magnitude of that parity if all occur. Require all target
differences of smaller magnitude in that class, and require the sum
neighbours at the chosen magnitude (none for the final choice).
\end{lemma}
\begin{proof}
Swapping the two coordinates of a target exchanges $d$ and $-d$ and
preserves $h$. Thus the required sums are the same whichever sign is
missing. A cover supplies an enumerated pair of magnitudes and contains
both required sets. For a diagonal family of deficit $\delta$, extend
its required set to every support of size $24-\delta$ in the integer
interval $[-L,L]$. Adjoin $\delta$ entries from this support in every
nondecreasing sequence, allowing repetition. This includes its actual
multiset, since the support uses one occurrence of each distinct value.

Test the quadratic identity above. For even axes both first moments
vanish. For odd axes the first moments $u,v$ must satisfy
$|u|+|v|\le22$ and $u\equiv v\bmod2$, by the formulas for $I,J$.
The odd-entry counts in the two families sum to 24 when $L=23$ and
agree when $L=22$, since $H-D=2j-L$. A placement would pass all
these tests. Exhaustion gives zero compatible keys for each of the six
even-axis and three odd-axis terminal branches.
\end{proof}

\begin{lemma}\label{lem:mixed-enum}
Once the exceptional queen is fixed, every cover in the mixed family
supplies a pair of complementary signatures in the enumeration below.
\end{lemma}
\begin{proof}
Remove all targets attacked by the exceptional queen and denote its
ranks by $d_0,h_0$. The remaining 23 queens have ranks in $[-22,22]$,
with $D,H$ of the same parity for each queen. Each remaining target has
ranks of opposite parity. For $b=0,1$ form the target graph with
difference vertices of parity $b$ and sum vertices of parity $1-b$.
Each body queen contributes exactly one entry to each graph.

In one graph enumerate the difference support by including or omitting
each available difference. Every omission forces all adjacent sum
vertices. A target with an endpoint outside the rank domain forces
the other endpoint, and one with both endpoints outside is impossible.
Extend the forced sums by every subset of the remaining sum vertices.
Give each selected vertex one entry, then distribute all remaining entries
among these vertices, with repetition, to make 23 entries in the graph.
This includes the support and multiplicity of every cover.
For each family count repetitions beyond the first occurrence and add one
if its exceptional rank also occurs in the body support.

Record the signature $(a,u,v,q,r,s)$ consisting of the number of
difference entries, the difference and sum first moments, the combined
square sum, and the two repetition counts just defined. Across both
graphs the signature must total $(23,0,0,4048,\delta_T,\delta_S)$.
Conversely the test only needs this necessary condition: matching
signatures could still fail to give a placement. The computed
intersection of one signature set with the complements of the other is
empty for each of the six mixed terminal branches.
\end{proof}

The original certificates encode the finite key comparisons as CNFs.
Independent code for the odd and mixed families starts from the omitted
\emph{sums} instead. For any actual cover let $F$ be the differences forced
by unavailable sums or adjacent to an omitted sum. Every member of $F$
is present. Every sum adjacent to a difference outside $F$ must be
present as well. Enumerating these unions and adjoining arbitrary extra
entries to each resulting subcover therefore includes the actual
multisets. Each terminal receipt identifies its original root and branch
by SHA-256. A separate test for even axes also gives a sharp bound.

\begin{corollary}\label{cor:twoholes}
Among permutations on the 24 even rows and 24 even columns of $Q_{47}$,
the minimum number of uncovered squares is two.
\end{corollary}
\begin{proof}
The $23^2$ odd--odd targets have $|d|+|h|\le22$ and $d\equiv h\bmod2$;
the inverse target coordinates are $(23+h+d,23+h-d)$.
Use the first missing magnitudes from Lemma~\ref{lem:pure-enum}. For a
missing difference of magnitude $a$, every sum of the same parity with
$|h|\le22-a$ is a neighbour. Each missing required sum leaves at least
one distinct target uncovered. Thus for at most one hole it suffices to
allow zero or one omission from the union of required sums. Complete
both lists to 24 entries from $[-23,23]$ in nondecreasing order of the
added entries, allowing arbitrary repetitions. Keep zero first moments,
combined square sum 4600 and combined odd count 24. The checker obtains
554 distinct ordered support pairs. Testing the targets for each pair
leaves none with at most one hole. For attainment, the permutation in
\file{data/frame/q47_24_two_holes.json} misses just $(1,5)$ and $(41,45)$
by direct inspection of the four line families at every square.
\end{proof}

The same mixed-family argument applies before the two occupied even
lines have been identified. For each possible count of occupied lines
of each parity consistent with a cube, let $r,s$ be the numbers of empty
rows and columns of a given parity pair. If both are positive, at least
$\lceil\max(r,s)/2\rceil$ queens of the corresponding square colour
are needed: a queen can cover at most two targets on a fixed empty line.
On 25 cubes across six frames, enumerating the row, column and queen
colour counts under these bounds forces all 23 odd lines and one even
line to be occupied on each axis. For each even row and column consistent
with the cube, we place the exceptional queen at their intersection and
apply the mixed-family moment test above. None of the 94 combinations of
intersection and diagonal deficits passes. These 25 exclusions cover
895 leaves of the earlier proof partition, including 34 SAT leaves.

\section{Computation and verification}
The initial SAT case is specified by the four exterior margins and the
four deficits. We enumerate frame rectangles up to the eight board
symmetries, then enumerate admissible deficit vectors as in \cite{Davies}.
At $n=47$, Proposition~\ref{prop:margin} removes 271,298 of 369,412 cases,
leaving 98,114. The parity and frame tests force a $12+12$ split in
65,512 of the retained cases. In the order used by the census, 40,807
have a nonsquare frame with odd longer span, 1,792 have an odd-span
square frame and $\delta_R+\delta_C>0$, and 22,913 additional cases
pass the nonorthodox deficit test ($\delta_R+\delta_C>2$, or $\delta_R+\delta_C=2$
with $\delta_R\ne\delta_C$). For example, in the nonorthodox $11+13$
case adding the two grid bounds gives $\delta_R+\delta_C\le2$. At
equality all four parity classes have even size by Lemma~\ref{lem:grid},
so $\delta_R=\delta_C=1$. Orthodoxy is possible only with margins
$(0,0),(1,1),(0,2)$ or $(2,0)$ on each axis: all empty indices must
have one parity, and at least 23 indices are empty. These tests together
with Theorem~\ref{thm:split} and Proposition~\ref{prop:profile} use only
the eight case parameters.

The baseline formula records queen locations and occupied lines. Each
square requires at least one of its four lines to be occupied. We impose
the prescribed line counts, the frame, and the perimeter bounds of
\cite{Davies}. A family of zero deficit has at most one queen per line.
Otherwise, when all four deficits are positive, a truncated unary counter
enforces exactly 24 queens. The baseline also uses symmetry breaking within
the case stabilizer, the square-frame bound on $E$ where applicable, and
the forced colour count where it can be read off a zero-deficit diagonal
family. For a difficult case we fix the occupation bits of the next inner
pair of rows and pair of columns. At depth $j$ the bits correspond in
order to rows $c+j,d-j$ and columns $a+j,b-j$, and a one means occupied.
All 16 assignments cover the parent whenever the two rows and two columns
are distinct.
Repeated splits give a tree of partial assignments, or \emph{cubes}.

For eligible square cases the solver can retry a timed-out cube with
the two-colour budgets of Section~\ref{sec:colour} before splitting
again. The extra clauses count wasted incidences by colour with truncated
totalizers. For a dominating placement, set each waste indicator to the
actual value of its geometric condition and select its actual budget.
The capacity identities then extend the placement to a satisfying
assignment. This implication is what permits a refutation of the
strengthened formula to close the original cube.

The paired timing experiment is reported in Appendix~\ref{app:pilot}.

\paragraph{Completion of the lower bound.}
The reconciled ledger closes all 98,114 retained roots. A root is covered
by a checked exclusion of the root itself or by exclusions covering every
assignment in its splits. A colour retry records digests for the baseline
cube and the strengthened formula. A profile exclusion links its shared
certificate to the line occupancies forced by the branch. For a reusable core the checker verifies the DRAT trace
against the core CNF itself. Each clause in that CNF must occur in the
hashed root or be a signed layer unit specified by the core mask.
The root's symmetry-breaking clauses are computed once from its
stabilizer before any splitting, so all its sibling cubes inherit them. The core therefore excludes every completion of its mask.
A leaf can use several masks if their union covers every completion of
the leaf, as checked recursively on the unfixed bits.

The partitions changed as arithmetic exclusions and symmetries were added.
The counts below refer to different stages of the same computation.
\begin{center}
\begin{tabular}{@{}ll@{}}
\toprule
Stage & Coverage \\
\midrule
Original search & 15 terminal branches across 13 frames (Section~\ref{sec:terminal})\\
Mixed-family reduction & 25 branches replace 895 former leaves in six frames\\
Published representatives & 17 mixed-moment leaves after whole-root symmetry\\
\bottomrule
\end{tabular}
\end{center}
Of the original 15 branches, five belong to roots covered by symmetry.
The other ten use four even-permutation, two odd-repeat and four mixed
moment exclusions.
Together with the 271,298 geometric exclusions this accounts for all
369,412 cases on canonical rectangles and rules out 24 queens. The placement in
Section~\ref{sec:bounds} gives the matching upper bound, proving
Theorem~\ref{thm:q47}.

\paragraph{Verification record.}
The closeout checker regenerates the canonical case census, reconciles
the original run ledgers and branch archives, replays the terminal
certificates and checks the 25-queen placement square by square. Some
early distributed runs discarded raw traces after successful DRAT checks.
For those cases the audit uses the verification receipts recorded at
solve time. For independent replay, retain the complete census and map each
eight-parameter case to its lexicographically least image under the eight
board symmetries. Transforming a placement transforms its extreme empty
lines and permutes its four deficits. Hence a dominating placement in a
referred case would give one in the representative. The representative's
own stabilizer can then put it in the orientation required by its SAT
formula. This argument applies to whole roots. The checker recomputes
each of the 7,097 references and requires its representative to have a
direct proof partition.

The 91,017 representatives have 131,547 leaves. Of these, 3,203 use
1,527 shared core certificates, 17 use mixed moments and 182 use
Proposition~\ref{prop:axis}. Of these, 171 have no surviving parity distribution,
five force the even permutation, three force the odd repeat and three
force the orthodox profile. The remaining 128,145 are SAT leaves.
In stored replay the index distinguishes 126,101 frame, 1,982 colour
and 62 historical formulas. The last group uses the earlier
\file{run_hard.py} encoder with symmetry breaking. In fresh replay those
62 use 16 current frame and 46 colour formulas. In all, fresh replay
uses 122,250 frame formulas and 5,895 colour formulas. It checks each
released formula digest, generates a DRAT proof and checks the proof
before discarding it. Instructions are in
\file{data/frame/q47_release/README.md}.

\paragraph{Implementation.}
At verifier revision \texttt{1ea84df} of \cite{software}, \file{code/frame/odd_margin_bounds.py}
reproduces the case counts and \file{odd_square_profile.py} adds the
reflected-line constraints. In the same directory,
\file{odd_colour_perimeter.py} builds the strengthened formula and
\file{q47_orthodox.py} rebuilds the shared certificate and checks a branch
receipt. The placement in Section~\ref{sec:bounds} is also recorded in
\file{data/frame/q47_24_one_hole.json}.
After installing the dependencies and pinned tools listed in
\file{data/frame/q47_release/README.md}, a referee can run the public
verification from the repository root:
\begin{verbatim}
.venv/bin/python code/frame/replay_q47_archive.py \
  data/frame/q47_release/proof-index.jsonl.zst \
  --regenerate --workers 8 --record build/q47-replay.jsonl
\end{verbatim}
This reconstructs the SAT formulas, checks fresh DRAT proofs, and checks
the applicability of each arithmetic and shared-core exclusion. The
record permits a restart after interruption. For the placements, census,
symmetry references and partitions alone, replace the final line with
\file{--index-only}. For those checks plus all 3,402 core and arithmetic
leaves, use \texttt{--certificates-only --workers 4}.
The separate program \file{code/frame/check_q47_terminal_enumeration.py}
checks the 15 original terminal branches. The historical ledger audit
with \file{code/frame/check_q47_closeout.py} additionally requires the
author's original ledgers and branch archives under \file{build/};
a clean checkout supplies the public replay above.

AI assistance was used in drafting the manuscript, implementing
the arithmetic reductions and verification software, and examining
the computational records. The release includes the executable checks,
the independent enumerations described above and their inputs.

\appendix
\section{Paired timing experiment}\label{app:pilot}
\begin{table}[htbp]
\centering\small
\begin{tabular}{llrrrr}
\toprule
Margins $(l,r,b,t)$ & Cube & \multicolumn{2}{c}{Baseline (s)}
 & \multicolumn{2}{c}{Colour budgets (s)}\\
\cmidrule(lr){3-4}\cmidrule(lr){5-6}
 & & Solve & Check & Solve & Check\\
\midrule
$(0,14,2,12)$ & \texttt{0110}         & cap   & ---  & 9.14  & 2.80\\
$(0,14,2,12)$ & \texttt{1111}         & cap   & ---  & 15.64 & 4.45\\
$(4,7,5,6)$   & \texttt{1010}         & cap   & ---  & 13.42 & 4.69\\
$(4,7,5,6)$   & \texttt{111110011001} & 14.55 & 4.62 & 2.00  & 0.55\\
\bottomrule
\end{tabular}
\caption{A paired local pilot on selected zero-deficit cubes. Each solve
had a 40 CPU-second limit. Check is CPU time for \texttt{drat-trim}.
A cap produced no proof to check.}\label{tab:pilot}
\end{table}

Table~\ref{tab:pilot} includes three previously timed-out cubes and one
solved control. Both encodings used the same CaDiCaL and checker binaries
on one pinned logical CPU. The 16 baseline child proofs for cube
\texttt{0110} totalled 59.71 MB and took 18.75 checker CPU-seconds in the search records. The
single strengthened proof was 9.30 MB and checked in 2.80 CPU-seconds.
On four whole frames in a separate ten-second pilot both encodings reached
the limit. The measurements quantify the benefit on these selected
branches. An overall speed estimate would require a representative sample.

\begin{thebibliography}{99}
\bibitem{Alon} N.~Alon, \emph{Combinatorial Nullstellensatz},
Combin.\ Probab.\ Comput.\ \textbf{8} (1999), 7--29.\newline
\href{https://doi.org/10.1017/S0963548398003411}{doi:10.1017/S0963548398003411}.
\bibitem{BM} A.~P.~Burger, C.~M.~Mynhardt, \emph{Properties of dominating
sets of the queens graph $Q_{4k+3}$}, Utilitas Math.\ \textbf{57} (2000),
237--253. \href{https://combinatorialpress.com/um/vol57/}{Journal archive}.
\bibitem{Weakley} W.~D.~Weakley, \emph{Domination in the queen's graph}, in:
Graph Theory, Combinatorics, and Algorithms, vol.~2,
Y.~Alavi and A.~Schwenk (eds.), Wiley, 1995, 1223--1232.
\bibitem{BGMW} S.~Boz\'oki, P.~G\'al, I.~Marosi, W.~D.~Weakley,
\emph{Domination of the rectangular queen's graph},
Electron.\ J.\ Combin.\ \textbf{26}(4) (2019), \#P4.45.
\href{https://doi.org/10.37236/6026}{doi:10.37236/6026}.
\bibitem{Chan} E.~Chan, \emph{Color-split frame reductions for queen domination
on even boards of sizes 50 through 64}, Hexagon, 2026.
\href{https://hexagonmath.org/pdf/2611.01169v1}{2611.01169v1}.
\href{https://github.com/CosmicCrusader23/queens/tree/0965e89b9ae745b1f96a86c6f8bbc509015684cb}{Computational records}.
\bibitem{Davies} J.~Davies, \emph{$\gamma(Q_{35})=19$ and certified queen
domination numbers up to $n=46$}, Hexagon, 2026.
\href{https://hexagonmath.org/2611.00830v1}{2611.00830v1}.
\bibitem{software} J.~Davies, \emph{Queen domination} (repository),
\href{https://github.com/jasondavies/queen-domination/tree/1ea84dfec92c2ed1584d2ca59550768a7dd34bd2}{verifier revision \texttt{1ea84df}}.
\bibitem{Djukic} A.~\DJ uki\'c, \emph{The Queen Domination Number of the
$48\times48$ Chessboard: $\gamma(Q_{48})=25$}, Zenodo, 2026.
\href{https://doi.org/10.5281/zenodo.24537290}{doi:10.5281/zenodo.24537290}.
\bibitem{FW} D.~Finozhenok, W.~D.~Weakley, \emph{An improved lower bound for
domination numbers of the queen's graph}, Australas.\ J.\ Combin.\
\textbf{37} (2007), 295--300.
\href{https://ajc.maths.uq.edu.au/pdf/37/ajc_v37_p295.pdf}{Journal PDF}.
\bibitem{HC} T.~Rostami, \emph{Hint and Conquer: New Queen Domination Numbers
by SAT Solving}, preprint, 2026.
\href{https://github.com/TahaRostami/Gamma/blob/hint-and-conquer/doc.pdf}{Manuscript}.
\bibitem{Neuhaus} S.~Neuhaus, \emph{Eigenschaften kleinster dominierender
Mengen und Dominanzzahlen von Damengraphen}, Dissertation, Universit\"at
zu K\"oln, 2009. \url{https://kups.ub.uni-koeln.de/2916/}.
\bibitem{OW} P.~R.~J.~\"Osterg\aa rd, W.~D.~Weakley, \emph{Values of domination
numbers of the queen's graph}, Electron.\ J.\ Combin.\ \textbf{8} (2001), \#R29.
\href{https://doi.org/10.37236/1573}{doi:10.37236/1573}.
\end{thebibliography}
\end{document}
